\documentclass[12pt]{article}
\usepackage{amsmath,amssymb,amsfonts,amsthm}
\usepackage[margin=1in]{geometry}

\begin{document}

\begin{center}
{\Large\bfseries Chapter 9, Problem 3}\\[6pt]
Byron \& Fuller, \textit{Mathematics of Classical and Quantum Physics}
\end{center}

\bigskip

\textbf{Problem.} Show that if $\{\lambda_n\}$ and $\{x_n(t)\}$ are respectively the eigenvalues and eigenvectors of a completely continuous Hermitian integral operator, then the sum
\[
S_N(t) = \sum_{n=1}^{N} (x_n, f)\,x_n(t)
\]
converges in the sense of mean convergence.

\bigskip
\textbf{Solution.}

By the completeness theorem (Theorem~9.4), the eigenvectors $\{x_n\}$ of a completely continuous Hermitian operator $A$ form a complete orthonormal set in the Hilbert space $H$. For any $f \in H$, the Fourier coefficients are $c_n = (x_n, f)$.

We need to show that the partial sums $S_N = \sum_{n=1}^{N} c_n\,x_n$ converge in the $L_2$ norm (mean convergence).

Consider the difference for $M > N$:
\[
\|S_M - S_N\|^2 = \left\|\sum_{n=N+1}^{M} c_n\,x_n\right\|^2 = \sum_{n=N+1}^{M} |c_n|^2,
\]
where we used the orthonormality of $\{x_n\}$.

By Bessel's inequality, we know that
\[
\sum_{n=1}^{\infty} |c_n|^2 \le \|f\|^2 < \infty.
\]

Since the series $\sum |c_n|^2$ converges, its tail $\sum_{n=N+1}^{M} |c_n|^2 \to 0$ as $N, M \to \infty$. This means $\{S_N\}$ is a Cauchy sequence in the Hilbert space $H$.

Since $H$ is complete (it is a Hilbert space), every Cauchy sequence converges. Therefore $S_N$ converges in the norm of $H$, i.e., in the sense of mean convergence:
\[
\left\|f - \sum_{n=1}^{N} (x_n, f)\,x_n\right\|^2 = \|f\|^2 - \sum_{n=1}^{N} |c_n|^2 \to 0 \quad \text{as } N \to \infty.
\]

The last equality uses the completeness relation (Parseval's equality), which holds because $\{x_n\}$ is a complete orthonormal set:
\[
\|f\|^2 = \sum_{n=1}^{\infty} |(x_n, f)|^2.
\]

Therefore $S_N \to f$ in the mean. \hfill $\blacksquare$

\end{document}
